% Numerical results for the two-dimensional stochastic optimization.
% Include from the main document with:  % Numerical results for the two-dimensional stochastic optimization.
% Include from the main document with:  % Numerical results for the two-dimensional stochastic optimization.
% Include from the main document with:  % Numerical results for the two-dimensional stochastic optimization.
% Include from the main document with:  \input{results_2d}
% Requires: booktabs, amsmath, siunitx (or plain tabular if siunitx is absent).
% NOTE: this file was authored without a local LaTeX toolchain; it has not been
% compiled.  Numbers are verbatim from runs 20261118-114831 and 20261120-111849
% (see DOCS/Results.md and DOCS/GPU.md).

\label{sec:results-2d}
\label{chap:results}

\section{Benchmark and certificate protocol}

The benchmark is a $200\times120$\,m aquifer discretized on $160\times96$ cells
($15\,360$ unknowns) with log-normal conductivity ($\sigma=0.65$, $\ell=15$\,m,
Karhunen--Lo\`eve expansion truncated at $95\%$ of the variance, $n_{KL}=1529$),
a $10\ \mathrm{m^3/d}$ source at $100\ \mathrm{mg/L}$ for $3650$\,d, and nine
candidate wells on a hexagonal lattice.  The chance constraint
\begin{equation}
  \mathbb{P}\bigl(G(\mathbf{q},\xi)\le 0\bigr)\ \ge\ 1-\alpha ,
  \qquad \alpha = 0.05 ,
\end{equation}
where $G$ is the worst compliance-point overshoot over the final $20\%$ of the
horizon, is handled by a conditional-value-at-risk penalty with Rockafellar--Uryasev
weights,
\begin{equation}
  \min_{\mathbf{q}\in\mathcal{Q}}\ \mathbb{E}[J(\mathbf{q},\xi)]
  + \rho_g\bigl[\mathrm{CVaR}_{\alpha}(G)\bigr]_+^2
  + \rho_h\bigl[\mathrm{CVaR}_{\alpha_h}(-H)\bigr]_+^2 ,
\end{equation}
evaluated on $M_{\text{train}}=100$ scenarios and searched by pattern search with
capacity projection and a batched compass sweep.

A policy is only reported as \emph{certified} after it has been re-evaluated on
fresh scenarios by an independent verifier, and the reliability is quoted with a
Wilson score interval,
\begin{equation}
  \frac{\hat p + \frac{z^2}{2n} \pm z\sqrt{\frac{\hat p(1-\hat p)}{n}+\frac{z^2}{4n^2}}}
       {1+\frac{z^2}{n}} ,
  \qquad z = 1.95996 ,
\end{equation}
because a bare fraction hides how much of the claim is sample size.

\section{Two certified policies}

\begin{table}[htbp]
\centering
\caption{Two independently certified optima of the same benchmark.  They differ
only in floating-point precision; both satisfy every gate.}
\label{tab:certified-policies}
\begin{tabular}{lrr}
\toprule
 & float64 & float32 (mixed) \\
\midrule
$\mathbf{q}^\star$ [$\mathrm{m^3/d}$] & $(0,9.72,5.86,0,3.38,0,0,0,0)$ & $(0,9.33,8.16,0,0,0,0,0,0)$ \\
total pumping [$\mathrm{m^3/d}$]      & $19.0$ & $17.5$ \\
active wells                          & $3$    & $2$ \\
$J$ (train / test)                    & $433\,945$ / $434\,008$ & $428\,592$ / $428\,579$ \\
$\mathrm{CVaR}_{5\%}(G)$ (train / test) & $+0.018$ / $-0.007$ & $+0.005$ / $+0.009$ \\
reliability (train / test)            & $0.990$ / $0.975$ & $0.990$ / $0.975$ \\
holdout during search                 & $1.000$ & $0.975$ \\
max.\ mass balance error              & $4.2\times10^{-11}$ & $2.3\times10^{-5}$ \\
wall clock (T4)                       & $4.14$\,h & $1.67$\,h \\
\bottomrule
\end{tabular}
\end{table}

Both optima lie \emph{on} the constraint, $\mathrm{CVaR}_\alpha(G)\approx 0$ from
either side, as the solution of a chance-constrained program must: they are the
cheapest policies that are still safe, not policies that happen to be safe.
Their costs agree to $1.2\%$ while their active well sets differ, so near the
optimum the objective is flat along the direction ``which wells carry the duty''.

\section{Out-of-sample certificate}

\begin{table}[htbp]
\centering
\caption{Reliability on 400 fresh scenarios with 95\% Wilson intervals.  Both
clear the $1-\alpha-1\%=0.94$ gate, but not equally.}
\label{tab:certificate}
\begin{tabular}{lrrr}
\toprule
policy & violations & reliability & 95\% Wilson CI \\
\midrule
float64 $\mathbf{q}^\star$ & $1/400$ & $0.998$ & $[0.986,\ 1.000]$ \\
float32 $\mathbf{q}^\star$ & $7/400$ & $0.983$ & $[0.964,\ 0.991]$ \\
\bottomrule
\end{tabular}
\end{table}

The two policies are indistinguishable at $M=40$, where both show $0$ violations
and a Wilson lower bound of $0.912$, and only separate at $M\ge200$.  Policy
ranking therefore requires an evaluation sample substantially larger than the
sample the search itself used --- a point that a single point-estimate
reliability cannot make.

Those \num{400} scenarios extend the run's training seed, so the first
\num{100} of them are the fields the optimizer was trained on.  Re-drawing
scenarios from a seed no run ever touched (\texttt{--seed-base 900000}) leaves the
certificate intact but changes the tail: $J = \num{428509}$ against the run's
\num{428592}, reliability $391/400 = 0.978$ with a Wilson interval of
$[0.958,\ 0.988]$, and $\mathrm{CVaR}_{5\%}(G) = +0.616$ instead of the $+0.005$
the run reports for its own training scenarios.  The tail statistic a run prints
is a statement about the sample it optimized on; only a fresh draw says what the
policy costs an aquifer.

\section{Which statistic can be certified}

A tail reported by a run is an average of the worst $\alpha$ fraction of that
run's own scenarios.  At $\alpha=0.05$ and $n=400$ that is $\approx20$ draws, and
they are the \emph{most variable} draws in the set, so the standard error of the
tail mean is $\operatorname{sd}(\text{tail})/\sqrt{20}$, not
$\operatorname{sd}(G)/\sqrt{n}$.  Re-solving five policies on one common draw of
\num{400} unseen scenarios makes the consequence concrete
(Table~\ref{tab:policy-ladder}).

\begin{table}[htbp]
\centering
\caption{Five policies re-solved on the same \num{400} unseen scenarios.  The
policy selected on only sixteen scenarios reported reliability $1.000$ on its own
sample and fails here; the two policies that pump six or nine wells are safe
beyond question and much more expensive.  The tail column carries its standard
error because it cannot be read to better than that.}
\label{tab:policy-ladder}
\begin{tabular}{lccccc}
\toprule
policy & $\mathbf{q}^\star$ wells & $J$ & $\mathrm{CVaR}_{5\%}(G)\pm$s.e. & reliability [95\% CI] & gate \\
\midrule
$M_{\text{train}}=16$            & $4$ & $\num{423041}$ & $+1.865\pm0.41$ & $0.912\ [0.881,0.936]$ & fail \\
accepted (float32)               & $2$ & $\num{428509}$ & $+0.616\pm0.47$ & $0.978\ [0.958,0.988]$ & pass \\
float64                          & $3$ & $\num{433924}$ & $+0.251\pm0.32$ & $0.980\ [0.961,0.990]$ & pass \\
six wells                        & $6$ & $\num{577917}$ & $-1.000\pm0.00$ & $1.000\ [0.990,1.000]$ & pass \\
nine wells                       & $9$ & $\num{821250}$ & $-1.000\pm0.00$ & $1.000\ [0.990,1.000]$ & pass \\
\bottomrule
\end{tabular}
\end{table}

Three readings follow.  First, the audit discriminates: a policy whose own
sixteen scenarios showed no violation fails decisively on \num{400} unseen ones,
so the certificate is a test, not a formality.  Second, safety costs money and the
first two percent costs the most: $+1.3\%$ of $J$ buys a smaller tail estimate,
while removing every observed violation costs $+35\%$ (six wells) or $+92\%$
(nine wells) --- and those two policies saturate the compliance metric at its
floor ($\bar G=-1$), so their $\mathrm{CVaR}=-1$ is a cap rather than a measured
margin.  Third, the tail does not order policies: $0.616$ against $0.251$ is not
a difference (see below), and $0.980$ against $0.978$ is not one either.

Re-solving the accepted policy on three disjoint draws separates the two
statistics (Table~\ref{tab:tail-draws}).

\begin{table}[htbp]
\centering
\caption{The same policy on three disjoint \num{400}-scenario draws.  The cost
reproduces to $0.02\%$ and the reliability to $0.003$; the tail moves by a factor
of $2.4$ and its spread is smaller than the standard error of any single value,
i.e.\ the three estimates are one estimate measured three times.}
\label{tab:tail-draws}
\begin{tabular}{rcrcc}
\toprule
seed base & $J$ & $\bar G$ & $\mathrm{CVaR}_{5\%}(G)$ & reliability [95\% CI] \\
\midrule
$900000$ & $\num{428509}$ & $-0.872$ & $+0.616$ ($\pm0.47$) & $0.978\ [0.958,0.988]$ \\
$950000$ & $\num{428605}$ & $-0.901$ & $+0.254$ ($\pm0.21$) & $0.975\ [0.955,0.986]$ \\
$990000$ & $\num{428579}$ & $-0.890$ & $+0.334$ ($\pm0.24$) & $0.978\ [0.958,0.988]$ \\
\bottomrule
\end{tabular}
\end{table}

The reporting rule adopted here is therefore: \emph{certify with reliability and
its Wilson interval}, because a binomial proportion estimated from all $\num{400}$
draws is stable, transportable across draws and backends, and comparable between
policies; \emph{quote $\mathrm{CVaR}$ with its standard error} and only as a
statement about where a policy sits relative to the constraint, never as a
validated number; and \emph{never rank policies by tail value} at the sample sizes
that are affordable.  The run's $+0.005$ and the fresh draw's $+0.62$ describe the
same policy and differ only in which twenty scenarios were drawn --- which is
precisely the sense in which an in-sample tail is not a validation.

\section{Effect on the plume}

Evaluated on $80$ held-out scenarios (medians), the float64 policy lowers the peak
concentration from $35.81$ to $0.433\ \mathrm{mg/L}$ and the compliance-point
concentration from $17.23$ to $0.0260\ \mathrm{mg/L}$, raising the number of
scenarios meeting the $1.0\ \mathrm{mg/L}$ limit from $1/80$ to $78/80$.  The
plume is displaced rather than destroyed: its peak moves down-gradient to
$x\approx115$\,m, adjacent to the active well.  For reference, no pumping gives
$J=322\,341$ with reliability $0.000$, and uniform half-capacity pumping gives
$J=1\,186\,250$ with reliability $1.000$.

\section{Computational cost}

The search is dominated by repeated transport solves and was run on a rented GPU.
Two measurements frame the cost: batching a whole compass sweep into one kernel
call is worth $47\times$ at $80\times40$ on a T4, and moving from float64 to a
mixed-precision float32 recurrence (float32 operator, float64 Krylov recurrence)
is worth $2.46\times$ at $160\times96$.  A third lever removes the per-iteration
kernel launches themselves: passing the stencil coefficients as arguments lets the
whole BiCGSTAB iteration be compiled into one call, which measured $2.2\times$ at
$80\times40$ and $2.8\times$ at $160\times96$ on an A100 with the objective
unchanged ($1.3\times10^{-7}$, float32).  The workload is latency-bound in float32
and fp64-throughput-bound in float64, so the choice of accelerator is
regime-dependent rather than monotone in price; see \texttt{DOCS/GPU.md} for the
audit and the honest ``do not rent a GPU for this'' cases.

% Requires: booktabs, amsmath, siunitx (or plain tabular if siunitx is absent).
% NOTE: this file was authored without a local LaTeX toolchain; it has not been
% compiled.  Numbers are verbatim from runs 20261118-114831 and 20261120-111849
% (see DOCS/Results.md and DOCS/GPU.md).

\label{sec:results-2d}
\label{chap:results}

\section{Benchmark and certificate protocol}

The benchmark is a $200\times120$\,m aquifer discretized on $160\times96$ cells
($15\,360$ unknowns) with log-normal conductivity ($\sigma=0.65$, $\ell=15$\,m,
Karhunen--Lo\`eve expansion truncated at $95\%$ of the variance, $n_{KL}=1529$),
a $10\ \mathrm{m^3/d}$ source at $100\ \mathrm{mg/L}$ for $3650$\,d, and nine
candidate wells on a hexagonal lattice.  The chance constraint
\begin{equation}
  \mathbb{P}\bigl(G(\mathbf{q},\xi)\le 0\bigr)\ \ge\ 1-\alpha ,
  \qquad \alpha = 0.05 ,
\end{equation}
where $G$ is the worst compliance-point overshoot over the final $20\%$ of the
horizon, is handled by a conditional-value-at-risk penalty with Rockafellar--Uryasev
weights,
\begin{equation}
  \min_{\mathbf{q}\in\mathcal{Q}}\ \mathbb{E}[J(\mathbf{q},\xi)]
  + \rho_g\bigl[\mathrm{CVaR}_{\alpha}(G)\bigr]_+^2
  + \rho_h\bigl[\mathrm{CVaR}_{\alpha_h}(-H)\bigr]_+^2 ,
\end{equation}
evaluated on $M_{\text{train}}=100$ scenarios and searched by pattern search with
capacity projection and a batched compass sweep.

A policy is only reported as \emph{certified} after it has been re-evaluated on
fresh scenarios by an independent verifier, and the reliability is quoted with a
Wilson score interval,
\begin{equation}
  \frac{\hat p + \frac{z^2}{2n} \pm z\sqrt{\frac{\hat p(1-\hat p)}{n}+\frac{z^2}{4n^2}}}
       {1+\frac{z^2}{n}} ,
  \qquad z = 1.95996 ,
\end{equation}
because a bare fraction hides how much of the claim is sample size.

\section{Two certified policies}

\begin{table}[htbp]
\centering
\caption{Two independently certified optima of the same benchmark.  They differ
only in floating-point precision; both satisfy every gate.}
\label{tab:certified-policies}
\begin{tabular}{lrr}
\toprule
 & float64 & float32 (mixed) \\
\midrule
$\mathbf{q}^\star$ [$\mathrm{m^3/d}$] & $(0,9.72,5.86,0,3.38,0,0,0,0)$ & $(0,9.33,8.16,0,0,0,0,0,0)$ \\
total pumping [$\mathrm{m^3/d}$]      & $19.0$ & $17.5$ \\
active wells                          & $3$    & $2$ \\
$J$ (train / test)                    & $433\,945$ / $434\,008$ & $428\,592$ / $428\,579$ \\
$\mathrm{CVaR}_{5\%}(G)$ (train / test) & $+0.018$ / $-0.007$ & $+0.005$ / $+0.009$ \\
reliability (train / test)            & $0.990$ / $0.975$ & $0.990$ / $0.975$ \\
holdout during search                 & $1.000$ & $0.975$ \\
max.\ mass balance error              & $4.2\times10^{-11}$ & $2.3\times10^{-5}$ \\
wall clock (T4)                       & $4.14$\,h & $1.67$\,h \\
\bottomrule
\end{tabular}
\end{table}

Both optima lie \emph{on} the constraint, $\mathrm{CVaR}_\alpha(G)\approx 0$ from
either side, as the solution of a chance-constrained program must: they are the
cheapest policies that are still safe, not policies that happen to be safe.
Their costs agree to $1.2\%$ while their active well sets differ, so near the
optimum the objective is flat along the direction ``which wells carry the duty''.

\section{Out-of-sample certificate}

\begin{table}[htbp]
\centering
\caption{Reliability on 400 fresh scenarios with 95\% Wilson intervals.  Both
clear the $1-\alpha-1\%=0.94$ gate, but not equally.}
\label{tab:certificate}
\begin{tabular}{lrrr}
\toprule
policy & violations & reliability & 95\% Wilson CI \\
\midrule
float64 $\mathbf{q}^\star$ & $1/400$ & $0.998$ & $[0.986,\ 1.000]$ \\
float32 $\mathbf{q}^\star$ & $7/400$ & $0.983$ & $[0.964,\ 0.991]$ \\
\bottomrule
\end{tabular}
\end{table}

The two policies are indistinguishable at $M=40$, where both show $0$ violations
and a Wilson lower bound of $0.912$, and only separate at $M\ge200$.  Policy
ranking therefore requires an evaluation sample substantially larger than the
sample the search itself used --- a point that a single point-estimate
reliability cannot make.

Those \num{400} scenarios extend the run's training seed, so the first
\num{100} of them are the fields the optimizer was trained on.  Re-drawing
scenarios from a seed no run ever touched (\texttt{--seed-base 900000}) leaves the
certificate intact but changes the tail: $J = \num{428509}$ against the run's
\num{428592}, reliability $391/400 = 0.978$ with a Wilson interval of
$[0.958,\ 0.988]$, and $\mathrm{CVaR}_{5\%}(G) = +0.616$ instead of the $+0.005$
the run reports for its own training scenarios.  The tail statistic a run prints
is a statement about the sample it optimized on; only a fresh draw says what the
policy costs an aquifer.

\section{Which statistic can be certified}

A tail reported by a run is an average of the worst $\alpha$ fraction of that
run's own scenarios.  At $\alpha=0.05$ and $n=400$ that is $\approx20$ draws, and
they are the \emph{most variable} draws in the set, so the standard error of the
tail mean is $\operatorname{sd}(\text{tail})/\sqrt{20}$, not
$\operatorname{sd}(G)/\sqrt{n}$.  Re-solving five policies on one common draw of
\num{400} unseen scenarios makes the consequence concrete
(Table~\ref{tab:policy-ladder}).

\begin{table}[htbp]
\centering
\caption{Five policies re-solved on the same \num{400} unseen scenarios.  The
policy selected on only sixteen scenarios reported reliability $1.000$ on its own
sample and fails here; the two policies that pump six or nine wells are safe
beyond question and much more expensive.  The tail column carries its standard
error because it cannot be read to better than that.}
\label{tab:policy-ladder}
\begin{tabular}{lccccc}
\toprule
policy & $\mathbf{q}^\star$ wells & $J$ & $\mathrm{CVaR}_{5\%}(G)\pm$s.e. & reliability [95\% CI] & gate \\
\midrule
$M_{\text{train}}=16$            & $4$ & $\num{423041}$ & $+1.865\pm0.41$ & $0.912\ [0.881,0.936]$ & fail \\
accepted (float32)               & $2$ & $\num{428509}$ & $+0.616\pm0.47$ & $0.978\ [0.958,0.988]$ & pass \\
float64                          & $3$ & $\num{433924}$ & $+0.251\pm0.32$ & $0.980\ [0.961,0.990]$ & pass \\
six wells                        & $6$ & $\num{577917}$ & $-1.000\pm0.00$ & $1.000\ [0.990,1.000]$ & pass \\
nine wells                       & $9$ & $\num{821250}$ & $-1.000\pm0.00$ & $1.000\ [0.990,1.000]$ & pass \\
\bottomrule
\end{tabular}
\end{table}

Three readings follow.  First, the audit discriminates: a policy whose own
sixteen scenarios showed no violation fails decisively on \num{400} unseen ones,
so the certificate is a test, not a formality.  Second, safety costs money and the
first two percent costs the most: $+1.3\%$ of $J$ buys a smaller tail estimate,
while removing every observed violation costs $+35\%$ (six wells) or $+92\%$
(nine wells) --- and those two policies saturate the compliance metric at its
floor ($\bar G=-1$), so their $\mathrm{CVaR}=-1$ is a cap rather than a measured
margin.  Third, the tail does not order policies: $0.616$ against $0.251$ is not
a difference (see below), and $0.980$ against $0.978$ is not one either.

Re-solving the accepted policy on three disjoint draws separates the two
statistics (Table~\ref{tab:tail-draws}).

\begin{table}[htbp]
\centering
\caption{The same policy on three disjoint \num{400}-scenario draws.  The cost
reproduces to $0.02\%$ and the reliability to $0.003$; the tail moves by a factor
of $2.4$ and its spread is smaller than the standard error of any single value,
i.e.\ the three estimates are one estimate measured three times.}
\label{tab:tail-draws}
\begin{tabular}{rcrcc}
\toprule
seed base & $J$ & $\bar G$ & $\mathrm{CVaR}_{5\%}(G)$ & reliability [95\% CI] \\
\midrule
$900000$ & $\num{428509}$ & $-0.872$ & $+0.616$ ($\pm0.47$) & $0.978\ [0.958,0.988]$ \\
$950000$ & $\num{428605}$ & $-0.901$ & $+0.254$ ($\pm0.21$) & $0.975\ [0.955,0.986]$ \\
$990000$ & $\num{428579}$ & $-0.890$ & $+0.334$ ($\pm0.24$) & $0.978\ [0.958,0.988]$ \\
\bottomrule
\end{tabular}
\end{table}

The reporting rule adopted here is therefore: \emph{certify with reliability and
its Wilson interval}, because a binomial proportion estimated from all $\num{400}$
draws is stable, transportable across draws and backends, and comparable between
policies; \emph{quote $\mathrm{CVaR}$ with its standard error} and only as a
statement about where a policy sits relative to the constraint, never as a
validated number; and \emph{never rank policies by tail value} at the sample sizes
that are affordable.  The run's $+0.005$ and the fresh draw's $+0.62$ describe the
same policy and differ only in which twenty scenarios were drawn --- which is
precisely the sense in which an in-sample tail is not a validation.

\section{Effect on the plume}

Evaluated on $80$ held-out scenarios (medians), the float64 policy lowers the peak
concentration from $35.81$ to $0.433\ \mathrm{mg/L}$ and the compliance-point
concentration from $17.23$ to $0.0260\ \mathrm{mg/L}$, raising the number of
scenarios meeting the $1.0\ \mathrm{mg/L}$ limit from $1/80$ to $78/80$.  The
plume is displaced rather than destroyed: its peak moves down-gradient to
$x\approx115$\,m, adjacent to the active well.  For reference, no pumping gives
$J=322\,341$ with reliability $0.000$, and uniform half-capacity pumping gives
$J=1\,186\,250$ with reliability $1.000$.

\section{Computational cost}

The search is dominated by repeated transport solves and was run on a rented GPU.
Two measurements frame the cost: batching a whole compass sweep into one kernel
call is worth $47\times$ at $80\times40$ on a T4, and moving from float64 to a
mixed-precision float32 recurrence (float32 operator, float64 Krylov recurrence)
is worth $2.46\times$ at $160\times96$.  A third lever removes the per-iteration
kernel launches themselves: passing the stencil coefficients as arguments lets the
whole BiCGSTAB iteration be compiled into one call, which measured $2.2\times$ at
$80\times40$ and $2.8\times$ at $160\times96$ on an A100 with the objective
unchanged ($1.3\times10^{-7}$, float32).  The workload is latency-bound in float32
and fp64-throughput-bound in float64, so the choice of accelerator is
regime-dependent rather than monotone in price; see \texttt{DOCS/GPU.md} for the
audit and the honest ``do not rent a GPU for this'' cases.

% Requires: booktabs, amsmath, siunitx (or plain tabular if siunitx is absent).
% NOTE: this file was authored without a local LaTeX toolchain; it has not been
% compiled.  Numbers are verbatim from runs 20261118-114831 and 20261120-111849
% (see DOCS/Results.md and DOCS/GPU.md).

\label{sec:results-2d}
\label{chap:results}

\section{Benchmark and certificate protocol}

The benchmark is a $200\times120$\,m aquifer discretized on $160\times96$ cells
($15\,360$ unknowns) with log-normal conductivity ($\sigma=0.65$, $\ell=15$\,m,
Karhunen--Lo\`eve expansion truncated at $95\%$ of the variance, $n_{KL}=1529$),
a $10\ \mathrm{m^3/d}$ source at $100\ \mathrm{mg/L}$ for $3650$\,d, and nine
candidate wells on a hexagonal lattice.  The chance constraint
\begin{equation}
  \mathbb{P}\bigl(G(\mathbf{q},\xi)\le 0\bigr)\ \ge\ 1-\alpha ,
  \qquad \alpha = 0.05 ,
\end{equation}
where $G$ is the worst compliance-point overshoot over the final $20\%$ of the
horizon, is handled by a conditional-value-at-risk penalty with Rockafellar--Uryasev
weights,
\begin{equation}
  \min_{\mathbf{q}\in\mathcal{Q}}\ \mathbb{E}[J(\mathbf{q},\xi)]
  + \rho_g\bigl[\mathrm{CVaR}_{\alpha}(G)\bigr]_+^2
  + \rho_h\bigl[\mathrm{CVaR}_{\alpha_h}(-H)\bigr]_+^2 ,
\end{equation}
evaluated on $M_{\text{train}}=100$ scenarios and searched by pattern search with
capacity projection and a batched compass sweep.

A policy is only reported as \emph{certified} after it has been re-evaluated on
fresh scenarios by an independent verifier, and the reliability is quoted with a
Wilson score interval,
\begin{equation}
  \frac{\hat p + \frac{z^2}{2n} \pm z\sqrt{\frac{\hat p(1-\hat p)}{n}+\frac{z^2}{4n^2}}}
       {1+\frac{z^2}{n}} ,
  \qquad z = 1.95996 ,
\end{equation}
because a bare fraction hides how much of the claim is sample size.

\section{Two certified policies}

\begin{table}[htbp]
\centering
\caption{Two independently certified optima of the same benchmark.  They differ
only in floating-point precision; both satisfy every gate.}
\label{tab:certified-policies}
\begin{tabular}{lrr}
\toprule
 & float64 & float32 (mixed) \\
\midrule
$\mathbf{q}^\star$ [$\mathrm{m^3/d}$] & $(0,9.72,5.86,0,3.38,0,0,0,0)$ & $(0,9.33,8.16,0,0,0,0,0,0)$ \\
total pumping [$\mathrm{m^3/d}$]      & $19.0$ & $17.5$ \\
active wells                          & $3$    & $2$ \\
$J$ (train / test)                    & $433\,945$ / $434\,008$ & $428\,592$ / $428\,579$ \\
$\mathrm{CVaR}_{5\%}(G)$ (train / test) & $+0.018$ / $-0.007$ & $+0.005$ / $+0.009$ \\
reliability (train / test)            & $0.990$ / $0.975$ & $0.990$ / $0.975$ \\
holdout during search                 & $1.000$ & $0.975$ \\
max.\ mass balance error              & $4.2\times10^{-11}$ & $2.3\times10^{-5}$ \\
wall clock (T4)                       & $4.14$\,h & $1.67$\,h \\
\bottomrule
\end{tabular}
\end{table}

Both optima lie \emph{on} the constraint, $\mathrm{CVaR}_\alpha(G)\approx 0$ from
either side, as the solution of a chance-constrained program must: they are the
cheapest policies that are still safe, not policies that happen to be safe.
Their costs agree to $1.2\%$ while their active well sets differ, so near the
optimum the objective is flat along the direction ``which wells carry the duty''.

\section{Out-of-sample certificate}

\begin{table}[htbp]
\centering
\caption{Reliability on 400 fresh scenarios with 95\% Wilson intervals.  Both
clear the $1-\alpha-1\%=0.94$ gate, but not equally.}
\label{tab:certificate}
\begin{tabular}{lrrr}
\toprule
policy & violations & reliability & 95\% Wilson CI \\
\midrule
float64 $\mathbf{q}^\star$ & $1/400$ & $0.998$ & $[0.986,\ 1.000]$ \\
float32 $\mathbf{q}^\star$ & $7/400$ & $0.983$ & $[0.964,\ 0.991]$ \\
\bottomrule
\end{tabular}
\end{table}

The two policies are indistinguishable at $M=40$, where both show $0$ violations
and a Wilson lower bound of $0.912$, and only separate at $M\ge200$.  Policy
ranking therefore requires an evaluation sample substantially larger than the
sample the search itself used --- a point that a single point-estimate
reliability cannot make.

Those \num{400} scenarios extend the run's training seed, so the first
\num{100} of them are the fields the optimizer was trained on.  Re-drawing
scenarios from a seed no run ever touched (\texttt{--seed-base 900000}) leaves the
certificate intact but changes the tail: $J = \num{428509}$ against the run's
\num{428592}, reliability $391/400 = 0.978$ with a Wilson interval of
$[0.958,\ 0.988]$, and $\mathrm{CVaR}_{5\%}(G) = +0.616$ instead of the $+0.005$
the run reports for its own training scenarios.  The tail statistic a run prints
is a statement about the sample it optimized on; only a fresh draw says what the
policy costs an aquifer.

\section{Which statistic can be certified}

A tail reported by a run is an average of the worst $\alpha$ fraction of that
run's own scenarios.  At $\alpha=0.05$ and $n=400$ that is $\approx20$ draws, and
they are the \emph{most variable} draws in the set, so the standard error of the
tail mean is $\operatorname{sd}(\text{tail})/\sqrt{20}$, not
$\operatorname{sd}(G)/\sqrt{n}$.  Re-solving five policies on one common draw of
\num{400} unseen scenarios makes the consequence concrete
(Table~\ref{tab:policy-ladder}).

\begin{table}[htbp]
\centering
\caption{Five policies re-solved on the same \num{400} unseen scenarios.  The
policy selected on only sixteen scenarios reported reliability $1.000$ on its own
sample and fails here; the two policies that pump six or nine wells are safe
beyond question and much more expensive.  The tail column carries its standard
error because it cannot be read to better than that.}
\label{tab:policy-ladder}
\begin{tabular}{lccccc}
\toprule
policy & $\mathbf{q}^\star$ wells & $J$ & $\mathrm{CVaR}_{5\%}(G)\pm$s.e. & reliability [95\% CI] & gate \\
\midrule
$M_{\text{train}}=16$            & $4$ & $\num{423041}$ & $+1.865\pm0.41$ & $0.912\ [0.881,0.936]$ & fail \\
accepted (float32)               & $2$ & $\num{428509}$ & $+0.616\pm0.47$ & $0.978\ [0.958,0.988]$ & pass \\
float64                          & $3$ & $\num{433924}$ & $+0.251\pm0.32$ & $0.980\ [0.961,0.990]$ & pass \\
six wells                        & $6$ & $\num{577917}$ & $-1.000\pm0.00$ & $1.000\ [0.990,1.000]$ & pass \\
nine wells                       & $9$ & $\num{821250}$ & $-1.000\pm0.00$ & $1.000\ [0.990,1.000]$ & pass \\
\bottomrule
\end{tabular}
\end{table}

Three readings follow.  First, the audit discriminates: a policy whose own
sixteen scenarios showed no violation fails decisively on \num{400} unseen ones,
so the certificate is a test, not a formality.  Second, safety costs money and the
first two percent costs the most: $+1.3\%$ of $J$ buys a smaller tail estimate,
while removing every observed violation costs $+35\%$ (six wells) or $+92\%$
(nine wells) --- and those two policies saturate the compliance metric at its
floor ($\bar G=-1$), so their $\mathrm{CVaR}=-1$ is a cap rather than a measured
margin.  Third, the tail does not order policies: $0.616$ against $0.251$ is not
a difference (see below), and $0.980$ against $0.978$ is not one either.

Re-solving the accepted policy on three disjoint draws separates the two
statistics (Table~\ref{tab:tail-draws}).

\begin{table}[htbp]
\centering
\caption{The same policy on three disjoint \num{400}-scenario draws.  The cost
reproduces to $0.02\%$ and the reliability to $0.003$; the tail moves by a factor
of $2.4$ and its spread is smaller than the standard error of any single value,
i.e.\ the three estimates are one estimate measured three times.}
\label{tab:tail-draws}
\begin{tabular}{rcrcc}
\toprule
seed base & $J$ & $\bar G$ & $\mathrm{CVaR}_{5\%}(G)$ & reliability [95\% CI] \\
\midrule
$900000$ & $\num{428509}$ & $-0.872$ & $+0.616$ ($\pm0.47$) & $0.978\ [0.958,0.988]$ \\
$950000$ & $\num{428605}$ & $-0.901$ & $+0.254$ ($\pm0.21$) & $0.975\ [0.955,0.986]$ \\
$990000$ & $\num{428579}$ & $-0.890$ & $+0.334$ ($\pm0.24$) & $0.978\ [0.958,0.988]$ \\
\bottomrule
\end{tabular}
\end{table}

The reporting rule adopted here is therefore: \emph{certify with reliability and
its Wilson interval}, because a binomial proportion estimated from all $\num{400}$
draws is stable, transportable across draws and backends, and comparable between
policies; \emph{quote $\mathrm{CVaR}$ with its standard error} and only as a
statement about where a policy sits relative to the constraint, never as a
validated number; and \emph{never rank policies by tail value} at the sample sizes
that are affordable.  The run's $+0.005$ and the fresh draw's $+0.62$ describe the
same policy and differ only in which twenty scenarios were drawn --- which is
precisely the sense in which an in-sample tail is not a validation.

\section{Effect on the plume}

Evaluated on $80$ held-out scenarios (medians), the float64 policy lowers the peak
concentration from $35.81$ to $0.433\ \mathrm{mg/L}$ and the compliance-point
concentration from $17.23$ to $0.0260\ \mathrm{mg/L}$, raising the number of
scenarios meeting the $1.0\ \mathrm{mg/L}$ limit from $1/80$ to $78/80$.  The
plume is displaced rather than destroyed: its peak moves down-gradient to
$x\approx115$\,m, adjacent to the active well.  For reference, no pumping gives
$J=322\,341$ with reliability $0.000$, and uniform half-capacity pumping gives
$J=1\,186\,250$ with reliability $1.000$.

\section{Computational cost}

The search is dominated by repeated transport solves and was run on a rented GPU.
Two measurements frame the cost: batching a whole compass sweep into one kernel
call is worth $47\times$ at $80\times40$ on a T4, and moving from float64 to a
mixed-precision float32 recurrence (float32 operator, float64 Krylov recurrence)
is worth $2.46\times$ at $160\times96$.  A third lever removes the per-iteration
kernel launches themselves: passing the stencil coefficients as arguments lets the
whole BiCGSTAB iteration be compiled into one call, which measured $2.2\times$ at
$80\times40$ and $2.8\times$ at $160\times96$ on an A100 with the objective
unchanged ($1.3\times10^{-7}$, float32).  The workload is latency-bound in float32
and fp64-throughput-bound in float64, so the choice of accelerator is
regime-dependent rather than monotone in price; see \texttt{DOCS/GPU.md} for the
audit and the honest ``do not rent a GPU for this'' cases.

% Requires: booktabs, amsmath, siunitx (or plain tabular if siunitx is absent).
% NOTE: this file was authored without a local LaTeX toolchain; it has not been
% compiled.  Numbers are verbatim from runs 20261118-114831 and 20261120-111849
% (see DOCS/Results.md and DOCS/GPU.md).

\label{sec:results-2d}
\label{chap:results}

\section{Benchmark and certificate protocol}

The benchmark is a $200\times120$\,m aquifer discretized on $160\times96$ cells
($15\,360$ unknowns) with log-normal conductivity ($\sigma=0.65$, $\ell=15$\,m,
Karhunen--Lo\`eve expansion truncated at $95\%$ of the variance, $n_{KL}=1529$),
a $10\ \mathrm{m^3/d}$ source at $100\ \mathrm{mg/L}$ for $3650$\,d, and nine
candidate wells on a hexagonal lattice.  The chance constraint
\begin{equation}
  \mathbb{P}\bigl(G(\mathbf{q},\xi)\le 0\bigr)\ \ge\ 1-\alpha ,
  \qquad \alpha = 0.05 ,
\end{equation}
where $G$ is the worst compliance-point overshoot over the final $20\%$ of the
horizon, is handled by a conditional-value-at-risk penalty with Rockafellar--Uryasev
weights,
\begin{equation}
  \min_{\mathbf{q}\in\mathcal{Q}}\ \mathbb{E}[J(\mathbf{q},\xi)]
  + \rho_g\bigl[\mathrm{CVaR}_{\alpha}(G)\bigr]_+^2
  + \rho_h\bigl[\mathrm{CVaR}_{\alpha_h}(-H)\bigr]_+^2 ,
\end{equation}
evaluated on $M_{\text{train}}=100$ scenarios and searched by pattern search with
capacity projection and a batched compass sweep.

A policy is only reported as \emph{certified} after it has been re-evaluated on
fresh scenarios by an independent verifier, and the reliability is quoted with a
Wilson score interval,
\begin{equation}
  \frac{\hat p + \frac{z^2}{2n} \pm z\sqrt{\frac{\hat p(1-\hat p)}{n}+\frac{z^2}{4n^2}}}
       {1+\frac{z^2}{n}} ,
  \qquad z = 1.95996 ,
\end{equation}
because a bare fraction hides how much of the claim is sample size.

\section{Two certified policies}

\begin{table}[htbp]
\centering
\caption{Two independently certified optima of the same benchmark.  They differ
only in floating-point precision; both satisfy every gate.}
\label{tab:certified-policies}
\begin{tabular}{lrr}
\toprule
 & float64 & float32 (mixed) \\
\midrule
$\mathbf{q}^\star$ [$\mathrm{m^3/d}$] & $(0,9.72,5.86,0,3.38,0,0,0,0)$ & $(0,9.33,8.16,0,0,0,0,0,0)$ \\
total pumping [$\mathrm{m^3/d}$]      & $19.0$ & $17.5$ \\
active wells                          & $3$    & $2$ \\
$J$ (train / test)                    & $433\,945$ / $434\,008$ & $428\,592$ / $428\,579$ \\
$\mathrm{CVaR}_{5\%}(G)$ (train / test) & $+0.018$ / $-0.007$ & $+0.005$ / $+0.009$ \\
reliability (train / test)            & $0.990$ / $0.975$ & $0.990$ / $0.975$ \\
holdout during search                 & $1.000$ & $0.975$ \\
max.\ mass balance error              & $4.2\times10^{-11}$ & $2.3\times10^{-5}$ \\
wall clock (T4)                       & $4.14$\,h & $1.67$\,h \\
\bottomrule
\end{tabular}
\end{table}

Both optima lie \emph{on} the constraint, $\mathrm{CVaR}_\alpha(G)\approx 0$ from
either side, as the solution of a chance-constrained program must: they are the
cheapest policies that are still safe, not policies that happen to be safe.
Their costs agree to $1.2\%$ while their active well sets differ, so near the
optimum the objective is flat along the direction ``which wells carry the duty''.

\section{Out-of-sample certificate}

\begin{table}[htbp]
\centering
\caption{Reliability on 400 fresh scenarios with 95\% Wilson intervals.  Both
clear the $1-\alpha-1\%=0.94$ gate, but not equally.}
\label{tab:certificate}
\begin{tabular}{lrrr}
\toprule
policy & violations & reliability & 95\% Wilson CI \\
\midrule
float64 $\mathbf{q}^\star$ & $1/400$ & $0.998$ & $[0.986,\ 1.000]$ \\
float32 $\mathbf{q}^\star$ & $7/400$ & $0.983$ & $[0.964,\ 0.991]$ \\
\bottomrule
\end{tabular}
\end{table}

The two policies are indistinguishable at $M=40$, where both show $0$ violations
and a Wilson lower bound of $0.912$, and only separate at $M\ge200$.  Policy
ranking therefore requires an evaluation sample substantially larger than the
sample the search itself used --- a point that a single point-estimate
reliability cannot make.

Those \num{400} scenarios extend the run's training seed, so the first
\num{100} of them are the fields the optimizer was trained on.  Re-drawing
scenarios from a seed no run ever touched (\texttt{--seed-base 900000}) leaves the
certificate intact but changes the tail: $J = \num{428509}$ against the run's
\num{428592}, reliability $391/400 = 0.978$ with a Wilson interval of
$[0.958,\ 0.988]$, and $\mathrm{CVaR}_{5\%}(G) = +0.616$ instead of the $+0.005$
the run reports for its own training scenarios.  The tail statistic a run prints
is a statement about the sample it optimized on; only a fresh draw says what the
policy costs an aquifer.

\section{Which statistic can be certified}

A tail reported by a run is an average of the worst $\alpha$ fraction of that
run's own scenarios.  At $\alpha=0.05$ and $n=400$ that is $\approx20$ draws, and
they are the \emph{most variable} draws in the set, so the standard error of the
tail mean is $\operatorname{sd}(\text{tail})/\sqrt{20}$, not
$\operatorname{sd}(G)/\sqrt{n}$.  Re-solving five policies on one common draw of
\num{400} unseen scenarios makes the consequence concrete
(Table~\ref{tab:policy-ladder}).

\begin{table}[htbp]
\centering
\caption{Five policies re-solved on the same \num{400} unseen scenarios.  The
policy selected on only sixteen scenarios reported reliability $1.000$ on its own
sample and fails here; the two policies that pump six or nine wells are safe
beyond question and much more expensive.  The tail column carries its standard
error because it cannot be read to better than that.}
\label{tab:policy-ladder}
\begin{tabular}{lccccc}
\toprule
policy & $\mathbf{q}^\star$ wells & $J$ & $\mathrm{CVaR}_{5\%}(G)\pm$s.e. & reliability [95\% CI] & gate \\
\midrule
$M_{\text{train}}=16$            & $4$ & $\num{423041}$ & $+1.865\pm0.41$ & $0.912\ [0.881,0.936]$ & fail \\
accepted (float32)               & $2$ & $\num{428509}$ & $+0.616\pm0.47$ & $0.978\ [0.958,0.988]$ & pass \\
float64                          & $3$ & $\num{433924}$ & $+0.251\pm0.32$ & $0.980\ [0.961,0.990]$ & pass \\
six wells                        & $6$ & $\num{577917}$ & $-1.000\pm0.00$ & $1.000\ [0.990,1.000]$ & pass \\
nine wells                       & $9$ & $\num{821250}$ & $-1.000\pm0.00$ & $1.000\ [0.990,1.000]$ & pass \\
\bottomrule
\end{tabular}
\end{table}

Three readings follow.  First, the audit discriminates: a policy whose own
sixteen scenarios showed no violation fails decisively on \num{400} unseen ones,
so the certificate is a test, not a formality.  Second, safety costs money and the
first two percent costs the most: $+1.3\%$ of $J$ buys a smaller tail estimate,
while removing every observed violation costs $+35\%$ (six wells) or $+92\%$
(nine wells) --- and those two policies saturate the compliance metric at its
floor ($\bar G=-1$), so their $\mathrm{CVaR}=-1$ is a cap rather than a measured
margin.  Third, the tail does not order policies: $0.616$ against $0.251$ is not
a difference (see below), and $0.980$ against $0.978$ is not one either.

Re-solving the accepted policy on three disjoint draws separates the two
statistics (Table~\ref{tab:tail-draws}).

\begin{table}[htbp]
\centering
\caption{The same policy on three disjoint \num{400}-scenario draws.  The cost
reproduces to $0.02\%$ and the reliability to $0.003$; the tail moves by a factor
of $2.4$ and its spread is smaller than the standard error of any single value,
i.e.\ the three estimates are one estimate measured three times.}
\label{tab:tail-draws}
\begin{tabular}{rcrcc}
\toprule
seed base & $J$ & $\bar G$ & $\mathrm{CVaR}_{5\%}(G)$ & reliability [95\% CI] \\
\midrule
$900000$ & $\num{428509}$ & $-0.872$ & $+0.616$ ($\pm0.47$) & $0.978\ [0.958,0.988]$ \\
$950000$ & $\num{428605}$ & $-0.901$ & $+0.254$ ($\pm0.21$) & $0.975\ [0.955,0.986]$ \\
$990000$ & $\num{428579}$ & $-0.890$ & $+0.334$ ($\pm0.24$) & $0.978\ [0.958,0.988]$ \\
\bottomrule
\end{tabular}
\end{table}

The reporting rule adopted here is therefore: \emph{certify with reliability and
its Wilson interval}, because a binomial proportion estimated from all $\num{400}$
draws is stable, transportable across draws and backends, and comparable between
policies; \emph{quote $\mathrm{CVaR}$ with its standard error} and only as a
statement about where a policy sits relative to the constraint, never as a
validated number; and \emph{never rank policies by tail value} at the sample sizes
that are affordable.  The run's $+0.005$ and the fresh draw's $+0.62$ describe the
same policy and differ only in which twenty scenarios were drawn --- which is
precisely the sense in which an in-sample tail is not a validation.

\section{Effect on the plume}

Evaluated on $80$ held-out scenarios (medians), the float64 policy lowers the peak
concentration from $35.81$ to $0.433\ \mathrm{mg/L}$ and the compliance-point
concentration from $17.23$ to $0.0260\ \mathrm{mg/L}$, raising the number of
scenarios meeting the $1.0\ \mathrm{mg/L}$ limit from $1/80$ to $78/80$.  The
plume is displaced rather than destroyed: its peak moves down-gradient to
$x\approx115$\,m, adjacent to the active well.  For reference, no pumping gives
$J=322\,341$ with reliability $0.000$, and uniform half-capacity pumping gives
$J=1\,186\,250$ with reliability $1.000$.

\section{Computational cost}

The search is dominated by repeated transport solves and was run on a rented GPU.
Two measurements frame the cost: batching a whole compass sweep into one kernel
call is worth $47\times$ at $80\times40$ on a T4, and moving from float64 to a
mixed-precision float32 recurrence (float32 operator, float64 Krylov recurrence)
is worth $2.46\times$ at $160\times96$.  A third lever removes the per-iteration
kernel launches themselves: passing the stencil coefficients as arguments lets the
whole BiCGSTAB iteration be compiled into one call, which measured $2.2\times$ at
$80\times40$ and $2.8\times$ at $160\times96$ on an A100 with the objective
unchanged ($1.3\times10^{-7}$, float32).  The workload is latency-bound in float32
and fp64-throughput-bound in float64, so the choice of accelerator is
regime-dependent rather than monotone in price; see \texttt{DOCS/GPU.md} for the
audit and the honest ``do not rent a GPU for this'' cases.
